% Hand-written; the numbers are quoted from Tables 1-3. Numbers change only through the repository.
\begin{table*}[t]
\centering
\scriptsize
\setlength{\tabcolsep}{3pt}
\begin{tabular}{p{3.4cm}p{6.2cm}p{5.6cm}}
\toprule
Report & Why & Provided by \\
\midrule
A company-disjoint split next to the temporal one &
On the time-ordered splits in use, 44--80\% of test calls come from companies in
training, and a probe reads the company out of a frozen call embedding 89.5\% of
the time (Tables~\ref{tab:substrate}, \ref{tab:controls}). &
\tickerdisjoint{} splits for FinCall-Surprise and MAEC (0.0\% overlap). The audit
command reports the share of test companies seen in training and scores seen and
unseen companies separately. \\
A gap of at least the longest horizon between training and test &
Five of six shipped boundaries have a zero-day gap while targets reach 30 days past
the call, so training labels contain test-period outcomes. &
Released temporal splits with a 30-session gap. The audit command reports the gap
in business days and counts training windows that reach into the test period. \\
Past-volatility baselines (persistence, HAR-RV, and a constant on change targets) &
They are hard to beat and often missing from published comparisons. On the measured
anchor the constant already scores 0.04--0.20 on $\Delta v$ (Table~\ref{tab:grid}). &
Stage~0 of the harness and the floors in the audit command, including a
company-only model. \\
Same-company and random swaps &
Where companies recur in training, a same-company swap moves the text model by
less than 0.03 while a random swap hurts, which attributes the score to the
company rather than the call (Table~\ref{tab:controls}). &
\texttt{ecvol controls} for feature-level swaps. Prediction-side swaps in the
audit command for any prediction file. \\
Day~0 measured from the release time &
Under the call-date anchor, day~0 moves for 53--72\% of calls and the reaction
session sits in the pre window (\S\ref{sec:timing}). &
Release times from 8-K filings in the timing tables. Call-date targets shipped as a
labeled variant for comparison with prior work. \\
The horizon definition, calendar or trading days &
The two definitions share only 0.78 correlation at $\tau{=}3$ and the persistence
MSE is 44--63\% higher under calendar days, so an MSE without its definition cannot
be compared. &
Targets under both definitions in every archive and a convention column in every
result table. \\
\bottomrule
\end{tabular}
\caption{\textbf{The reporting checklist of \S\ref{sec:recommendations}.} Each item
with the evidence behind it and the part of the release that provides it.}
\label{tab:checklist}
\end{table*}
